% Auto-generated by Exercise3.jl — do not edit by hand.
\begin{table}[htbp]
  \centering
  \caption{Structural Change and Aggregate Labor Productivity Growth}
  \label{tab:ex3_channels_growth}
  \begin{threeparttable}
  \centering
  \begin{tabular}{l S[table-format=-1.3] S[table-format=-1.3] S[table-format=-1.3]}
    \toprule
     & {1950--2023} & {1950--1980} & {1980--2023} \\
    \midrule
    \multicolumn{4}{l}{\textit{Panel A. Annualized growth rates}} \\
    Fixed-shares benchmark & 2.136 & 2.279 & 2.037 \\
    Baseline model & 2.040 & 2.095 & 2.001 \\
    Income effects removed & 2.091 & 2.041 & 2.126 \\
    Full shutdown & 2.280 & 2.063 & 2.432 \\
    \addlinespace
    \multicolumn{4}{l}{\textit{Panel B. Contributions to baseline growth}} \\
    I.\ Income effects & -0.051 & 0.055 & -0.126 \\
    II.\ Relative-price effects & -0.189 & -0.022 & -0.305 \\
    \midrule
    Total & -0.240 & 0.033 & -0.431 \\
    \bottomrule
  \end{tabular}
  \begin{tablenotes}[flushleft]
    \footnotesize
    \RaggedRight
    \item[] \textit{Notes:} Panel A reports annualized log growth rates of aggregate value added
  per production worker, $VA/L^p$, in percentage points. The fixed-shares benchmark weights
  baseline sectoral labor productivity growth by 1950 value-added shares. Income effects are
  removed by setting $\eta_j=1$; the full shutdown also sets $\sigma=1$. Each row of
  Panel A is a solved equilibrium. Panel B reports baseline growth minus growth with income effects removed
  (I) and growth with income effects removed minus full-shutdown growth (II); a negative entry
  means that the mechanism lowers aggregate labor productivity growth. Entries are rounded so
  that they add to the printed totals.
  \end{tablenotes}
  \end{threeparttable}
\end{table}
